
\documentclass[11pt]{article}
\usepackage[margin=1in]{geometry}
\usepackage{amsmath,amssymb}
\usepackage{enumitem}
\usepackage[T1]{fontenc}
\usepackage{lmodern}
\setlength{\parskip}{0.5em}
\setlength{\parindent}{0pt}

\begin{document}

\begin{center}
{\LARGE \textbf{IB Physics 2016--2020}}\\
{\large \textbf{Mapped Collection for D4}}\\
\vspace{0.3em}
\textit{Near-literal paraphrases matched to induction, generation, transmission, and capacitance.}
\end{center}

\textbf{Use.} These versions are designed to preserve the original solving experience: same scenario, same numerical data, and essentially the same progression of subparts.

\section*{Problems}

\begin{enumerate}[leftmargin=*, itemsep=1.1em]

\item \textbf{M17 TZ2 HL Paper 2, Q2(b) --- buoy driving a generator}\\
\textit{Diagram.} A buoy floats in a vertical tube in the sea. As the buoy rises, a cable turns a generator fixed on the sea bed. As the buoy falls, the cable rewinds but the generator does not deliver power. The buoy moves vertically with the water and is treated as performing SHM. The water wave has amplitude $4.3\,\text{m}$, wavelength $35\,\text{m}$, and speed $3.4\,\text{m s}^{-1}$.
\begin{enumerate}[label=(\alph*)]
\item Determine the maximum vertical speed of the buoy.
\item On a set of axes, sketch the generator output power against time for one complete oscillation of the buoy.
\item Explain why the generator produces power during only part of each cycle.
\end{enumerate}

\item \textbf{M17 TZ2 HL Paper 2, Q6(b)--(e) --- long-distance ac transmission}\\
\textit{Diagrams.} One figure shows a cable made from $32$ identical copper wires, each wire having length $35\,\text{km}$ and resistance $64\,\Omega$. Another figure shows a generator connected to a distant load through transmission cables. A later figure shows two identical cables operating in parallel. The generator output power is $110\,\text{MW}$ and the peak potential difference is $150\,\text{kV}$. The resistivity of copper is $1.7\times10^{-8}\,\Omega\text{m}$.
\begin{enumerate}[label=(\alph*)]
\item Calculate the radius of each copper wire.
\item Determine the peak current in the cable.
\item Calculate the power dissipated in the cable per unit length.
\item When two identical cables are used in parallel, determine the rms current in each cable.
\item Describe how the magnetic force between the two parallel cables changes throughout one ac cycle.
\item Explain why a step-up transformer is advantageous for transmitting electrical power over long distances.
\end{enumerate}

\item \textbf{M18 TZ2 HL Paper 2, Q7 --- pumped-storage hydroelectric station}\\
\textit{Diagram.} Water is stored behind a dam in an upper lake and then flows down through turbines to a lower lake. The vertical drop is $110\,\text{m}$. The station handles water at a rate of $1.2\times10^{5}\,\text{m}^3$ per minute. Take the density of water as $1.0\times10^{3}\,\text{kg m}^{-3}$.
\begin{enumerate}[label=(\alph*)]
\item Estimate the loss of gravitational potential energy per unit mass of water, including an appropriate unit.
\item Show that the average rate at which gravitational potential energy decreases is about $2.5\,\text{GW}$.
\item Explain how this type of station can still be economically worthwhile even though electrical energy is used to pump the water back up.
\end{enumerate}

\item \textbf{M19 TZ1 HL Paper 2, Q8(a)--(c) --- homemade capacitor and induced emf on discharge}\\
\textit{Diagrams.} One sketch shows a capacitor made from two aluminium-foil sheets with a sheet of plastic film between them; all three strips are $450\,\text{mm}$ wide. The plastic thickness is $55\,\mu\text{m}$ and its permittivity is $2.5\times10^{-11}\,\text{C}^{2}\text{N}^{-1}\text{m}^{-2}$. The capacitance is $68\,\text{nF}$. A second sketch shows the capacitor in a charge-discharge circuit with a $12\,\text{V}$ battery, a resistor of $1.2\,\text{k}\Omega$, a switch, and later a coil replacing the ammeter in the discharge branch.
\begin{enumerate}[label=(\alph*)]
\item Calculate the total length of aluminium foil needed.
\item The plastic breaks down if the electric field inside it exceeds $1.5\,\text{MN C}^{-1}$. Determine the maximum charge that can be stored safely.
\item When the fully charged capacitor is discharged through the resistor, calculate the time for the charge to fall to $50\%$ of its initial value.
\item The ammeter is replaced by a coil. Explain why an emf is induced in the coil while the capacitor discharges.
\item Suggest one change to the discharge circuit, other than altering the coil itself, that would increase the maximum induced emf.
\end{enumerate}

\item \textbf{N19 HL Paper 2, Q9 --- charge redistribution between capacitors}\\
\textit{Diagram.} Capacitor $X$ of capacitance $18\,\mu\text{F}$ is initially charged so that one plate carries $48\,\mu\text{C}$. Capacitor $Y$ of capacitance $12\,\mu\text{F}$ is initially uncharged. The two are then connected by a switch and resistor and allowed to reach a common potential difference.
\begin{enumerate}[label=(\alph*)]
\item Calculate the initial energy stored in capacitor $X$.
\item Determine the final charge on each capacitor after equilibrium is reached.
\item Find the total final electrical energy stored by the capacitor system.
\item Explain why there is a decrease in stored electrical energy when the capacitors are connected together.
\end{enumerate}

\item \textbf{N20 HL Paper 2, Q9 --- rotating coil and transmission}\\
\textit{Diagram.} A coil rotates uniformly in a magnetic field and provides an alternating emf to a transmission system. A later stage uses a step-up transformer before power is sent to a distant load. One graph relates induced emf to time.
\begin{enumerate}[label=(\alph*)]
\item Using magnetic flux linkage and Faraday's law, explain why the rotating coil produces an alternating emf.
\item If the peak output voltage is $25\times10^{3}\,\text{V}$, determine the rms voltage.
\item If the output power is $8.5\times10^{5}\,\text{W}$, calculate the rms current supplied.
\item Determine the required step-up factor if the transmission voltage is increased by a factor of $16$.
\item The peak emf is doubled and the rotation period is halved. State how the emf-time graph changes.
\end{enumerate}

\end{enumerate}

\end{document}
